\cotefile{1}
\scanpage{137}
Soient maintenant \inserted{$0<p,q,r\leq\infty$} $A\in L^p$ $A\in L^q$, $\frac1p+\frac1q\geq\frac1r$ démontrons la formule (6). Nous pouvons supposer \deleted{\unsure{$p,$}} que $p$ et $q$ \deleted{\illegible{2}}\inserted{sont finis} \deleted{\unsure{puisque deux infinis} (\unsure{on en aura alors}} \inserted{puisque si l'un était infini, les formules (2)} donneraient le résultat\unsure{;|.} \deleted{$\|AB\|_{\unsure{r}}$ ($\|AB\|_\infty=\|A\|_\infty\|B\|_\infty$) \unsure{est un}} \inserted{de plus,} peut supposer évidemment \deleted{\unsure{ici}} que l'on a exactement $\frac1p+\frac1q=\frac1r$ (puisque $\|X\|_r$ est fonction décroissante de $r$)

Soit $m$ \deleted{un} entier tel que $p\leq 2m$, $q\leq 2m$. \deleted{\unsure{Pour|Alors} $A$} Posons comme dans le lemme 2

\deleted{On a}

\[
\deleted{S_r(AB)\unsure{\leq|=} S_{\frac r2}(AB(AB)^*)=S_{\frac r{2m}}\unsure{A|(L)}}
\]

$L=(AB)^*(AB)^{\deleted{\unsure{*}}}$, $H=A^*A$, $K=B^*B$, on aura

\[
S_r(AB)=S_{\frac r2}(L)=S_{\frac r{2m}}(L^m)\deleted{\unsure{=r/2m}}
\]

d'où, en vertu de la formule (9), en posant pour la commodité $\lambda(x)=\frac{\pi}{\sin\pi x}$

\[
S_r(AB)=S_{\frac r{2m}}(L^m)=\frac{r/2m}{\lambda(r/2m)}
\int_0^\infty\frac{\log M_{L^m}(\rho)}{\rho^{1+r/2m}}\,d\rho \tag*{(11)}
\]

\deleted{\unsure{Or, on a} $\unsure{M_{L^m}|M}$} Or, $M_{L^m}(\rho)=\sum\alpha_n(L^m)\rho^n$, et \deleted{$M_{L^m}(\rho)=\sum_n\alpha_{\unsure{n|k}}(L^{\unsure{m|k}})\rho^{\unsure{k|n}}$} \inserted{$\alpha_n(L^m)\leq\alpha_n(H^m)\alpha_n(K^m)$} en vertu du lemme 2 on a \deleted{$\alpha_n(L^m)\leq\alpha_n(H^{\unsure{2|m}})\unsure{\alpha_n K}$}

% slanted
\margin{\uline{\unsure{Ici lemme|Ici le cours}} sur \uline{les fon}ctions \uline{entiè}res}

d'où

\[
M_{L^m}(\rho)\leq\sum\alpha_n(H^m)\alpha_n(K^m)\rho^n
\]

Posons $\left[p'=\frac1p\;\;q'=\frac1q\;\;r'=\frac1r\right]$, on a : \deleted{$p'+q'=r'$}

\[
\inserted{p'r+q'r=1}
\]

d'où

\[
\begin{aligned}
M_{L^m}(\rho)&\leq\sum\alpha_n(H^m)\alpha_n(K^m)\rho^{np'r+nq'r}\\
&=\sum_n\bigl(\alpha_n(H^m)\rho^{np'r}\bigr)
\bigl(\alpha_n(K^m)\rho^{n{\unsure{q|p}}'r}\bigr)
\end{aligned}
\]
% archive sheet number: 138
\scanpage{138}
\[
\begin{aligned}
&\leq\left(\sum_n\alpha_n(H^m)\rho^{np'r}\right)
\left(\sum_n\alpha_n(K^m)\rho^{nq'r}\right)\\
&\deleted{\unsure{=|\leq}}\qquad=M_{H^m}(\rho^{p'r})M_H(\rho^{q'r})
\end{aligned}
\]

d'où, compte tenu de (12)

\[
\frac{\lambda(r/2m)}{r/2m}S_r(AB)\leq
\int_0^\infty\frac{\log M_{H^m}(\rho^{p'r})}{\rho^{1+r/2m}}\,d\rho+
\int_0^\infty\frac{\log M_{K^m}(\rho^{q'r})}{\rho^{1+r/2m}}\,d\unsure{\rho}
\tag*{(13)}
\]

Calculons \deleted{\unsure{par exemple}} la première intégrale du second membre, par le changement de variable $\sigma=\rho^{p'r}$ d'où $\rho=\sigma^{pr'}$ $\frac{d\rho}{\rho}=pr'\frac{d\sigma}{\sigma}$ il vient

\[
\int_0^\infty\frac{\log M_{H^m}(\rho^{p'r})}{\rho^{r/2m}}\frac{d\rho}{\rho}
=pr'\int_0^\infty\frac{\log M_{H^m}(\sigma)}{\sigma^{p/2m}}\frac{d\sigma}{\sigma}
\]

Le deuxième membre est égal (\deleted{formule} (9)) à

\[
pr'\frac{\lambda(p/2m)}{p/2m}S_{p/2m}(H^m)=\frac{\lambda(\unsure{p/2m})}{r/2m}S_p(A).
\]

\unsure{De} $=$ la deuxième intégrale dans (13) est égale à $\frac{\lambda(q/2m)}{r/2m}S_q(B)$, et \deleted{on obtient} \inserted{(13) s'écrit ainsi}

\[
\lambda(r/2m)S_r(AB)\leq\lambda(p/2m)S_p(A)+\lambda(q/2m)S_q(B)
\]

ou aussi

\[
S_r(AB)\leq\frac{\sin\pi\frac{r}{2m}}{\sin\pi\frac{p}{2m}}S_p(A)+
\frac{\sin\pi\frac{r}{2m}}{\sin\pi\frac{q}{2m}}S_q(B)
\]

Cela prouve déjà que $S_{\unsure{r|p}}(AB)\leq+\infty$, et donne même une majoration effective de $S_r(AB)$ en fonction de $S_p(A)$ et $S_q(B)$. Faisant tendre $m$ vers $+\infty$, on obtient la formule plus simple

\[
\boxed{S_r(AB)\leq\frac{r}{\deleted{\unsure{2m}}p}S_p(A)+\frac rq S_q(B)}\tag*{(14)}
\]

C'est une inégalité remarquable, mais n'est pas celle que nous cherchons. Mais \deleted{cette relation \illegible{2}}
% archive sheet number: 139
\scanpage{139}
\deleted{\unsure{du fait de l'homogénéité|à part d'homogénéité} \unsure{il vient,|évidente,} qui permet}

Remplaçons dans cette formule $A$ par $\lambda A$, $B$ par $\mu B$ ($\lambda$ et $\mu$ nombres $\unsure{>| \geq}0$) on obtient :

\[
\deleted{\unsure{\lambda^p\mu^q=\frac{\lambda^p\mu^q}{r}S_r(AB)}}
\inserted{\frac{(\lambda\mu)^r}{r}S_r(AB)}
\leq\frac{\lambda^p}{p}S_p(A)+\frac{\mu^q}{q}S_q(B)
\]

\deleted{Exprimons que cette inégalité doit être vérifiée quels que soient $x,y\geq 0$.}

% slanted
\margin{\unsure{Ici lemme que} la condition du \unsure{n°} implique $c^{1/r}\leq a^{1/p}b^{1/q}$ i.e \uline{$c\leq a^\alpha b^\beta$}}

Pour simplifier, posons $a=S_p(A)$, $b=S_q(B)$, $c=S_{\unsure{r|p}}(C)$, on aura

\[
c\leq\mathop{\mathrm{Inf}}\limits_{\substack{\lambda\geq0\\\mu\geq0}}
\frac{a\frac{\lambda^p}{p}+b\frac{\mu^q}{q}}
{\deleted{\unsure{\frac{\lambda^p\mu^q}{r}}}\frac{(\lambda\mu)^r}{r}}
\]

\deleted{Posons $\unsure{x^p=\lambda,y^q}$}\deleted{$x=\lambda^p,\ y=\mu^q$,}

Posons $x=\lambda^p$, $y=\mu^q$, $\alpha=\frac rp$, $\beta=\frac rq$ d'où \unsure{,} \uline{$\frac rp+\frac rq=1$},

On a

\[
c\leq\mathop{\mathrm{Inf}}\limits_{\substack{x\unsure{>|\geq}0\\y>0}}
\frac{\alpha a\deleted{\unsure{\frac rp}}x+\deleted{\unsure{\frac rq}}\beta by}{x^\alpha y^\beta}
=\mathop{\mathrm{Inf}}\limits_{\substack{x\geq0\\y\geq0\\x^\alpha y^\beta=1}}\alpha a x+\beta by
\]

Or on vérifie aussitôt que la \deleted{minimum}\inserted{borne inférieure} du second membre est atteinte (car l'en\unsure{s.} des points $(x,y)$ tels que $x\geq0,y\geq0$ $ax+by\leq M$ -- valeur donnée -- est compact) elle est donc atteinte en un point $(x,y)$ tel que les équations \deleted{$a\,dx+b\,dy=0$}\inserted{$d(\alpha ax+\beta by)=0$}, $d(x^\alpha y^\beta)=0$ \deleted{soient compatibles} \inserted{aient une solution non nulle en} $dx,dy$  ;

or $d(x^\alpha y^\beta)=\alpha x^{\alpha-1}y^\beta\,dx+\beta y^{\beta-1}x^\alpha\,dy$, \deleted{i.e. tel que $\unsure{\det(\alpha a\quad\beta b\quad\alpha a}$} la condition d'extremum est donc $\frac{\alpha a}{\alpha x^{\alpha-1}y}=\frac{\beta b}{\beta x^\alpha y^{\beta-1}}$,

soit $ax=by$ \deleted{\unsure{$=$}}, \deleted{\unsure{et comme on a}} \inserted{d'où $ax=by=\frac{\alpha ax+\beta by}{\alpha+\beta}$} ; \deleted{$\unsure{ax+by=\frac{\alpha ax}{\alpha}=\frac{\alpha ax+\beta by}{\alpha}}$} Comme $\alpha+\beta=1$, le minimum cherché sera alors \deleted{\unsure{$ax+by$|$\alpha x+b$}} égal à $ax=by$ ; compte tenu de la condi%
% archive sheet number: 140
\scanpage{140}%
tion $x^\alpha y^\beta=1$ on trouve $x=\left(\frac ba\right)^\beta$, d'où \deleted{\unsure{nous|avons}} le minimum cherché $ax=b^\beta\deleted{\unsure{/a^\beta}}a^{1-\beta}=a^\alpha b^\beta$.

On trouve donc $c<a^\alpha b^\beta$, ce qui achève la démonstration.
% archive sheet number: 141

\scanpage{142}
\uline{Proposition} Pour que la suite $(H_i)$ \unsure{d'él. hermitiens|d'éch. hermitiens|d'hermitiens} $\in L^p(\mathscr H)$ converge vers $H_0\in L^p(\mathscr H)$, il faut et il suffit que

\[
\text{1°)}\quad\|H_i\|_p\to\|H_0\|_p\qquad
\text{(ou seulement }\limsup_{i\to\infty}\|H_i\|_p\leq\|H_0\|_p
\]

2°) Pour \deleted{tout} $\lambda$ du spectre de $H_0$, (soit $\alpha>0$ tel qu'il \deleted{$[\lambda-\alpha,\lambda+\alpha]$ ne \unsure{contien|contienne}}\inserted{n'existe pas de} valeur spectrale $\mu\neq\lambda$ de $H_0$ telle que $|\lambda\deleted{\unsure{+}}-\mu|\leq\unsure{\zeta|\alpha}$, \deleted{\unsure{élément spectral de $H$|l'ensemble spectral de $H$} \unsure{faut que}} \deleted{et telle} et tel que $\lambda-\alpha$ et $\lambda+\alpha$ n'appartiennent au spectre \deleted{de $H_{\unsure{0|i}}$}\inserted{d'aucun $H_i$}) \deleted{alors} la composante spectrale de $H_i$ \deleted{\unsure{relative}} dans $[\lambda-\alpha,\lambda+\alpha]$ est de dimension \deleted{\unsure{uniforme égale à}} qui finit par être égale à celle de $H_0$, et $h_i\to h_0$ (inutile de spécifier la topologie, en vertu d'une proposition préliminaire sur l'identité de \deleted{toutes} ces topologies)

% slanted
\margin{\unsure{évident} \unsure{faux ds} le cas non hermitien}

\uline{Corollaire 1} Les deux conditions suivantes sont néc. et suffis. \unsure{stes} pour la convergence de la suite $H_i$ vers $H_0$

\[
\text{1)}\quad\|H_i\|_p\to\|H_0\|_p\qquad\text{(resp$^{\mathrm t}$ \unsure{\ldots{} --})}
\]

\[
\text{2)}\quad H_i\to H_0\quad\text{dans }L^\infty(\mathscr H)
\]

\uline{Corollaire 2} $A\to |A|^r$ est application continue de $L^p$ dans $L^{\frac pr}$.

En effet, elle est composée de $A\to A^*A$, continue de $L^p$ dans $L^{\frac p2}$, avec
% archive sheet number: 143
\scanpage{143}
$H\to H^{\frac r2}$ de $L^{+\;\frac p2}$ dans $L^{\frac pr}$ ; or, les conditions de la prop. 1 sont aussitôt vérifiées \deleted{\unsure{en effet}} (condition 1 triviale ; condition 2 se voit \deleted{ainsi} en se ramenant à \deleted{\unsure{montrer que si}}\inserted{la dimension} finie, où c'est trivial ($\sum_{i=1}^{\unsure{n}}\lambda_i a_i\otimes a_i$, $\lambda_i\geq0$ et $a_i$ variables, $\to\lambda_i^{(0)}$ et $a_i^{(0)}$ ; on forme $\sum\lambda_i^{r/2}a_i\otimes a_i$ \unsure{\ldots{} --})

\uline{Remarque} 1 Si $0<p\leq2$, pour qu'une suite d'hermitiens $H_i$ dans $L^p$ converge \inserted{vers $H_0$}, il faut et il suffit qu'elle converge faiblement, et que $\|K_i\|\to\|H_0\|$ (ceci est connu au cas hilbertien)

% slanted
\margin{\deleted{\uline{\unsure{prendre comme|probablement}} \uline{\unsure{indépendant de $p$|pertinent ds $L^p$}} ! \unsure{la norme uniforme de Murray|On connaît le contre-exemple de Murray}}}

Est-ce vrai si $p<+\infty$ ? Uniforme convexité de $L^p$ ?

\uline{Remarque 2} Est-ce que $S_p\lvert H^r-K^r\rvert^{\frac1r}$ est fonction croissante de $r$ ($H$ hermitien positif) ? -- Il suffirait de savoir que $r<1$ implique

\[
\bigl\|\lvert H^r-K^r\rvert^{\frac1r}\bigr\|\leq\|H-K\|
\]

--- Alors on aurait une précision notable dans le corollaire 1.
% archive sheet number: 144

\scanpage{144}
\[
\left\{\begin{aligned}
\rho_{m+n-1}(vu)&\unsure{=|\leq}\rho_m(\unsure{v|u})\rho_n(u)\\
\rho_{m+n-1}(u+v)&\leq\rho_m(u)+\rho_n(v)
\end{aligned}\right.
\]

En effet,

\[
\rho_{m+n-1}(vu)=\mathop{\mathrm{Inf}}_{E_{m+n-1}}
\mathop{\mathrm{Sup}}_{\ni x,\,\|x\|\leq1}\|vu x\|
\]

\diagram{\begin{tikzcd}[ampersand replacement=\&,column sep=small]
E\arrow[r,"u"]\&F\arrow[r,"v"]\&\unsure{E|G}
\end{tikzcd}}

Or il existe $E_{n-1}\subset E$ tel que $\mathop{\mathrm{Sup}}_{\substack{\|x\|\leq1\\x\in E_{n-1}}}\deleted{\unsure{\rho}}\|ux\|=\rho_n(u)$ \deleted{\unsure{alors|et}} il existe $F_{m-1}\subset F$ tel que $\mathop{\mathrm{Sup}}_{\substack{\|y\|\leq1\\y\in F_{m-1}}}\|vy\|=\rho_m(v)$ soit alors \deleted{$\unsure{E_{m+1}|E_{m+n-2}}$}$H=u^{-1}(F_{m-1})$, on a

% margin note written sideways
% A tall left brace groups the two inequalities and their proofs.
\margin{\uline{majorations individuelles}}

\[
\mathop{\mathrm{Sup}}_{\substack{x\in H\cap E_{n-1}\\\|x\|\leq1}}\|vux\|\leq\rho_n(u)\rho_m(v)
\]

or $p-\deleted{\unsure{p}}1=\operatorname{codim}H\cap E_{n-1}\geq\operatorname{cod}H+\operatorname{cod}E_{n-1}\geq m-1+n-1=(m+n-1)-1$

donc

\[
\rho_p(vu)=\mathop{\mathrm{Inf}}\mathop{\mathrm{Sup}}_{\substack{x\in E_{p-1}\\\|x\|\leq1}}\|vux\|\leq\rho_n(u)\rho_m(v)
\]

à fortiori

\[
\boxed{\rho_{m+n-1}(vu)\leq\rho_n(u)\mathbin{.}\rho_m(v)}
\]

De $=$, si $u,v\in L_0(E,F)$, on a

\[
p-\deleted{\unsure{p}}1=\operatorname{codim}(E_{n-1}\unsure{,|\cap} E'_{m-1})\geq m-1+n-1=(m+n-1)\unsure{,|-1}
\]

\[
\deleted{\unsure{\rho_p(u+v)=\mathrm I|\rho_p(u+v)\leq}}
\mathop{\unsure{\mathrm{Inf}}}\mathop{\mathrm{Sup}}_{\substack{x\in E_{n-1}\cap E'_{m-1}\\\|x\|\leq1}}
\|(u+v)x\|\leq\rho_n(u)+\rho_m(v)
\]

d'où à fortiori

\[
\rho_{m+n-1}(u+v)\leq\rho_p(u+v)\leq\rho_n(u)+\rho_m(v)
\]

\deleted{\unsure{Cas|Est-ce}} On peut avoir $\limsup\frac{|\lambda_i(u)|}{\rho_i(u)}=+\infty$

Majorations individuelles pour $|\lambda_i(u)|\leq\rho_{\unsure{i/2|i/k}}(u)$ ou analogue ?
% archive sheet number: 145

\scanpage{145}
Si $E$ est un Banach, et $u_\alpha\to u$ dans $L_0(E)$, alors $\sigma(u_\alpha)\to\sigma(u)$ dans le sens suivant :

De plus, convergence des projecteurs spectraux \unsure{\ldots{} .} Dans le cas des op. hermitiens dans l'espa\unsure{ce} de Hilbert, on a $=$ là des conditions nécessaires et suffisantes de convergence.

Convergence de \unsure{l'hermitien|d'hermitiens|l'hermiticité} \inserted{$u_\alpha\to u$} dans $L^p(E)$ : il est nécessaire et suffisant que \textcircled{\scriptsize 1°} $\|u_\alpha\|_p\to\|u\|_p$, \textcircled{\scriptsize 2°} $u_\alpha\to u$ spectralement, i.e. dans $L(E)$

(il suffit \unsure{même} que l'on ait la convergence spectrale individuelle en chaque $\lambda\in\sigma(u)$, $\lambda\neq0$)

\uline{Corollaire 1} $\unsure{U|u}\to|A|$, de $L^p(E,F)$ dans $L^p(E)$ est continue. Plus généralement $u\to|u|^r$ de $L^p(E,F)$ dans $L^{p/r}(E)$ est continue.

Car composée de $u\to u^*u$ de $L^p(E,F)$ dans $(L^{p/2}(E))^+$, avec $v\to v^{\frac r2}$ de $(L^{p/2}(E))^+$ dans $L^{\frac pr}(E)$ ; or \deleted{\unsure{la|cette}} \unsure{dernière|dernier} est continue, car elle est continue de $L_0(E)^+$ dans $L_0(E)$ (ce dernier \unsure{point|fait} se vérifie \uline{spectralement} ; c'est sans doute vrai aussi de $L(E)^+$ dans $L(E)$ ? Inégalités précises ?)

\uline{Corollaire 2} Si $0<p\leq\unsure{2|1}$, pour qu'une suite \inserted{d'hermitiens} $u_\alpha\to u$ converge dans $L^p(E\deleted{\unsure{,F}})$, il faut et il suffit que $\|u_\alpha\|_p\to\|u\|_p$ et qu'on ait la convergence faible.
% archive sheet number: 146

\scanpage{146}
Question : \unsure{Sont uniformément convexes|Sont-ce des convexes|Sont-ils convexes|Dans quel cas (uniforme convexité} \unsure{les|ces|des} \unsure{espaces|esp.} $L^p(E,F)$ ? pour $0<p<\infty$ ? Pas d'importance.

2) Pour la continuité de $h\to h^r$, on est ramené au cas $0\leq r<1$ (car c'est trivial pour $r$ entier) Il semble possible que

\[
\bigl\|\lvert {\unsure{h|H}}^r-k^r\rvert\bigr\|^{1/r}\leq\|h-k\|\qquad\text{pour }0\leq r<1.
\]

Il \unsure{revient au $=$ de prouver|revient au $=$ de présumer} que

\[
S_{p/r}(h^r-k^r)=S_p\bigl(\lvert h^r-k^r\rvert^{1/r}\bigr)
\]

est fonction croissante de $r$, ($h$ et $k$ hermitiens positifs, $p$ quelconque).

On peut faire aussi $p=1$.
% archive sheet number: 147

\scanpage{147}
\uline{Th} Si $A_i\to A$ dans $L^p(\mathscr H)$, alors $\widehat A_{\unsure{i}}\to\widehat A$ dans $\Sigma_p$. Cela résulte facilement du

\uline{lemme} Pour tout $\varepsilon>0$, existe $\rho>0$ et $i_0$ tel que $i>i_0$ implique

\[
\sum_{r_\nu(A_i)\leq\rho}\bigl(r_\nu\unsure{(A_i)|(\lvert A_i\rvert)|[A_i]}\bigr)^p\leq\varepsilon
\]

Pour s'en convaincre, soit $\unsure{w|\omega}(t)$ la fonction convexe qui est identique à $e^{pt}$ pour $t\leq\log\rho$, et linéaire pour $t\geq\log\rho$, i.e $w(t)=p\rho^p(t-\log\rho)+\rho^p=\rho^p\log e\left(\frac{e^t}{\rho}\right)^p$ on a (posant $r_\nu=r_\nu(A_i)$, \deleted{\unsure{$s_\nu$ analogue}}\inserted{$s_\nu=r_\nu(|A_i|)$})

% margin note written sideways
\margin{\unsure{On n'a pas besoin} de la continuité de $A\to|A|$, si on remplace $|A|$ par $A^*A$}

\[
\sum w(\log r_\nu)\leq\sum w(\log(s_\nu))
\]

soit

\[
\sum_{r_\nu\leq\rho}r_\nu^p+\rho^p\sum_{r_\nu\geq\rho}\log e\left(\frac{r_\nu}{\rho}\right)^p
\leq\sum_{s_\nu\leq\rho}s_\nu^p+\rho^p\sum_{s_\nu\geq\rho}\log e\left(\frac{s_\nu}{\rho}\right)^p
\]

On peut donc \uline{supposer les $A_i\geq0$}, (car $|A_i|\to|A|$) et prouver que alors, pour tout $\varepsilon>0$, existe $\rho>0$ et $i_0$ tel que $i>i_0$ implique

\[
\sum_{r_\nu(A_i)\leq\rho}\unsure{r_\nu(A_i)^p|r_\nu\lvert A_i\rvert^p}
+\rho^p\sum_{r_\nu(A_i)\geq\rho}\log e\left(\frac{r_\nu(A_i)}{\rho}\right)^p<\varepsilon
\]

On prouve que d'abord, par extraction d'une \deleted{sous-suite $r_\nu\unsure{=}\,r_\nu(A_{\unsure{i'|i}})$, il suffit de montrer séparément}

\margin{\textcircled{\scriptsize \unsure{Si !|Si, Si}}}

\[
\deleted{\text{a)}\qquad\rho^p\sum_{r_\nu\geq\rho}\log e\left(\frac{r_\nu}{\rho}\right)^p\to0\qquad\text{si }\rho\to0}
\]

\deleted{b) $\rho$\unsure{$^{\prime}$} \unsure{énoncé du lemme|l'inégalité du lemme} pour les hermitiens positifs $A_i$.} \deleted{\unsure{Or} \illegible{7}}

\uline{Preuve de a)} Or, on a en effet

\[
\begin{aligned}
\rho^p\sum_{r_\nu\geq\rho}\log e\left(\frac{r_\nu}{\rho}\right)^p
&=\rho^p\int_{r>\rho}^{+\infty}-\loge\left(\frac{r}{\rho}\right)^p\,dn(r)\\
&=\rho^p\left[-n(r)\log\unsure{-|\ldots}\right]_{\rho+0}^{+\infty}
+p\rho^p\int_\rho^{+\infty}\frac{n(r)}{r}\,dr\\
&=\rho^p n(\rho)+p\rho^p\int_\rho^{+\infty}\frac{n(r)}{r}\,dr
\end{aligned}
\]

et il est bien connu, puisque $\sum r_\nu^p<\infty$, que $\rho^p n(\rho)\to0$
% archive sheet number: 148
\scanpage{148}
et que $\int_0^{+\infty}\frac{n(r)}{r^{1-p}}\,dr<+\infty$, à fortiori $\int_{\unsure{1|0|\rho}}^{+\infty}\frac{n(r)}r\,dr<+\infty$,

donc $p\rho^\alpha\int_\rho^{+\infty}\frac{n(r)}r\,dr<p\rho^\alpha\int_0^{+\infty}\frac{n(r)}r\,dr\to0$ si $\rho\to0$.

\deleted{pour tout $\varepsilon$, on peut supposer \unsure{aussi} que}

\[
\deleted{\sum\|A_{\deleted{\unsure{i}}\unsure{i+1}}-A_i\|_p<+\infty,\quad
\sum S_p(A_{i+1}-A_{i})<+\infty}
\]

\deleted{alors $A=\sum A_{i+1}-A_i$, donc $A\leq\sum|A_{i+1}-A_i|$}

\uline{Preuve de b} On a

\[
\begin{aligned}
\sum_{r_\nu(A_i)\leq\rho}(r_\nu(A_i))^{\unsure{r|p}}&=S_p(A_i)-\sum_{r_\nu(A_i)>\rho}(r_\nu(A_i))^p\\
&\longrightarrow\deleted{\unsure{\sum|\leq}} S_p(A)-\sum_{r_\nu>\rho}r_\nu^p=\sum_{r_\nu\leq\rho}r_\nu^p
\end{aligned}
\]

il suffit de prendre $\rho$ tel que $\sum_{r_\nu\leq\rho}r_\nu^p<\varepsilon$, puis $i$ assez grand.
% archive sheet number: 149

\scanpage{149}
$\mathscr H$ Hilbert. Si $u$ est à trace, $\frac1{1+u}=1-u\frac{R(u)}{\det(1+u)}$

\[
R(u)=\sum_0^\infty P_n(u)=\frac{\det(1+u)}{1+u}
\]

\uline{Théorème} \uline{1) Soit $u$ opérateur \deleted{\unsure{herm.}} à trace normal. Alors}

\[
\boxed{|P_n(u)|\leq P_n(|u|)\ ;\ |R(u)|\leq R(|u|)}\tag*{\unsure{$(*),\,R)$|$(*)$|$(*),\,({**})$}}
\]

En particulier, si $u$ est $\geq0$, il en est de $=$ de $P_n(u)$ et de $R(u)$

2) Si $u\geq0$, alors la suite des termes $\alpha_{n-k}(u)u^k$ dans l'expression de $P_n(u)$, $P_n(u)=\sum_{k=0}^n\alpha_{n-k}(u)u^k$, est une suite \uline{décroissante} d'hermitiens positifs \deleted{\unsure{quand $k$}}.

Donc on a en particulier, pour $u\geq0$

\[
\deleted{\unsure{\alpha_n(u)}}\quad
\boxed{0\leq P_n(u)\leq\alpha_n(u)\mathbin{.}1\qquad\text{d'où}\quad
\deleted{\unsure{R(u)\leq\alpha_n(u)\,1}}\inserted{0\leq R(u)\leq\det(1+u)1}}\tag*{(**)}
\]

\deleted{\unsure{3)}} \uline{Corollaire} Si $u$ est opérateur \deleted{\unsure{de rang}} à trace normal, on a

\[
\left\{\begin{aligned}
\|R(u)\|_\infty&\leq\|R(|u|)\|_\infty\leq\det(1+|u|)\leq e^{\|u\|_1}\\
\|P_n(u)\|_\infty&\leq\|P_n(|u|)\|_\infty\leq\alpha_n(|u|)\leq\frac{\|u\|_1^n}{n!}
\end{aligned}\right.\tag*{(***)}
\]

\uline{Démonstration} On démontre le

\uline{Lemme} Si \inserted{$u$ est opérateur normal} $u=\sum_1^\infty\lambda_i e_i\otimes e'_i$, alors $P_n(u)=\sum\unsure{p_i(u)|\rho_i(u)}e_i\otimes e_i$, où $p_i(u)=\sum_{(I)}'\lambda_{i_1}\cdots\lambda_{i_n}$, $\sum'$ étant étendue à tous les \deleted{\unsure{systèmes}} \deleted{ensembles}\inserted{$I=\{i_1\ldots i_n\}$} de $n$ indices, \deleted{\unsure{qui tels que $I$} \illegible{2}}\inserted{ne contenant pas $i$}

Pour le voir, on calcule successivement (dans la base $(e_i)$) \deleted{\unsure{$u^n$}}$u^n-$$\alpha_1(u)u^{n-1}$, $u^n-\alpha_1(u)u^{n-1}+\alpha_2(u)u^{n-1}$ par récurrence on obtient $\sum_{k=0}^{\unsure{p|n}}u^{n-k}\alpha_k(u)(-1)^k=(-1)^n\sum'\lambda_{i_1}\cdots\lambda_{i_k}$ d'où \deleted{\unsure{$P_n=$ pour tout}} le résultat.

Maintenant 1) devient évident. 2) Se vérifie sur les valeurs propres de façon immédiate.

\uline{Remarque} Notons que si $\rho_i$ sont les valeurs propres de $R(u)$, on a $\rho_i=\prod_{j\neq i}(1+\lambda_j)$, d'où $|\rho_i|\leq\prod_j(1+|\lambda_j|)=\det(1+\lvert u\rvert)$, on retrouve une
% archive sheet number: 150
\scanpage{150}
partie de $***$. Mais si \deleted{\unsure{$u\geq0$}}\inserted{$\unsure{1|-1}$ n'est pas valeur propre,} on a aussi, \deleted{pour tout} \inserted{\unsure{pour la dimension infinie|particulièrement en dimension infinie}} $\rho_i\deleted{\unsure{\geq}}=\frac{\det(1+\unsure{\lambda|u})}{1+\lambda_i}$ \deleted{\unsure{$\geq$|$\leq$}} d'où

\[
\boxed{\|R(u)\|_\infty\geq|\det(1+u)|}\quad\text{(\unsure{op normal|si dim infinie})}
\]

(car il existe des $\lambda_i$ arbitrairement petits, donc des $\rho_i$ arbitrairement \deleted{distants \unsure{par}}\inserted{voisins} de $\det(1+u)$ !

On \deleted{\unsure{en}} prouve donc que l'inégalité $***$ a) ne peut s'améliorer, car elle devient égalité pour $u\geq0$ !
% archive sheet number: 151
